\documentclass[10pt]{article}

\title{A Bibliography of Ægypt by John Crowley}
\section{Introduction}\label{intro}

\emph{Compiled by Brian Murphy}

The goal is to list all of the books and authors mentioned throughout
\emph{Ægypt}, or \emph{The Solitudes}, by chapter. Footnotes to each
chapter heading will list the authors that are mentioned within that
chapter. And, a complete list of those authors are listed at the end in
alphabetical order by the first name. If a chapter is missing below,
then it likely did not have a book that was mentioned within it.

\subsection[Author's Note]{\texorpdfstring{Author's
Note\footnote{Authors Mentioned: Joseph Campbell, Elizabeth L.
  Eisenstein, Mircea Eliade, Peter French, Hans Jonas, Frank E. \& Fritz
  P. Manuel, Giorgio di Santillana, Stephen Schoenbaum, Wayne Shumaker,
  Keith Thomas, Lynn Thorndike, D. P. Walker, Dame Francis Yates, John
  Mitchell, Katerine Maltwood, Robert Graves, Lois Rose, Richard Deacon,
  Walter Scott, Luis de Góngora.}}{Author's Note}}\label{authors-note43}

\begin{itemize}
\tightlist
\item
  Walter Scott. \emph{Hermetica}.
\item
  Luis de Góngora, trans. Cunningham, Gilbert P. \emph{Soledades}.
\end{itemize}

\subsection{The Prologue on Earth}\label{the-prologue-on-earth}

\begin{itemize}
\tightlist
\item
  Bruno, Giordano. \emph{De monade, numero et figura}. 1961.
\end{itemize}

\subsection[Vita - One]{\texorpdfstring{Vita -
One\footnote{Authors Mentioned: Peter Ramus}}{Vita - One}}\label{vita---one39}

\begin{itemize}
\tightlist
\item
  De Góngora Y Argote, Luis. \emph{The Solitudes of Luis De Góngora}.
  1968.
\end{itemize}

\subsection{Vita - Two}\label{vita---two}

\begin{itemize}
\tightlist
\item
  Kraft, Fellowes. \emph{Darkling Plain}.\footnote{This is a fictional
    book created by John Crowley.}
\item
  Burnett, Frances Hodgson, and Sue Ullstein. \emph{The Secret Garden}.
  Longman Publishing Group, 1975.
\item
  Norton, Mary. \emph{The Borrowers Abroad}.\footnote{This is a
    fictional book created by John Crowley.}\footnote{There is no novel
    by Norton titled \emph{The Borrowers Abroad}. This may have been
    deliberate by Crowley, or an accident, since all of the books are
    named in a similar pattern (i.e.~\emph{The Borrowers Afield},
    \emph{The Borrowers Afloat}, \emph{The Borrowers Afield}, \emph{The
    Borrowers Avenged}. More research could be done to see if this was a
    proposed title of another book in the series, or if this was a
    subtle way to shift the book away from the reality, or history in
    this case, that we experience as the reader.)}
\item
  Kraft, Fellowes. \emph{Under Saturn, A Novel of
  Wallenstein}\footnote{This is a fictional book created by John
    Crowley.}\footnote{Seems to be based on the biography of
    Wallenstein: Watson, Francis. \emph{Wallenstein, Soldier Under
    Saturn}. 1969.}
\item
  Kraft, Fellowes. \emph{Bitten Apples}\footnote{This is a fictional
    book created by John Crowley.}
\end{itemize}

\subsection[Vita - Three]{\texorpdfstring{Vita -
Three\footnote{Authors Mentioned: Lady Mary Wortley Montagu, Luis de
  Góngora, Percy Bysshe Shelley}}{Vita - Three}}\label{vita---three40}

\begin{itemize}
\tightlist
\item
  De Góngora Y Argote, Luis. \emph{The Solitudes of Luis De Góngora}.
  1968.\footnote{Crowley quotes Cunningham's translation of the First
    Solitude lines 94-96 and then 108-111, then 112-116. Then switches
    to the original Spanish for pieces of lines 113-114, and then again
    the Spainish for lines 119-121}
\end{itemize}

\subsection{Vita - Four}\label{vita---four}

\begin{itemize}
\tightlist
\item
  Barr, Frank Walker. \emph{Time's Body}.\footnote{This is a fictional
    book created by John Crowley.}
\item
  Barr, Frank Walker. \emph{Mythos and Tyrannos}.\footnote{This is a
    fictional book created by John Crowley.}
\item
  \emph{Seven Sobs of a Sorrowful Soule for Sinne}. 1597.
\item
  Shakespeare, William. \emph{Measure for Measure}. Cambridge UP, 2006.
\item
  \emph{Unknown Volumes of Encyclopedia Britannica}. 1939.\footnote{I am
    assuming the 1939 edition since Pierce mentions bringing several
    volumes of this fictitious version with him to Blackberry Jams.}
\item
  Malory, Thomas, and Sidney Lanier. \emph{The Boy's King Arthur: Being
  Sir Thomas Malory's History of King Arthur and His Knights of the
  Round Table}. 1921.
\item
  Kraft, Fellowes. \emph{Bruno's Journey}. 1931\footnote{This is a
    fictional book created by John Crowley.}
\end{itemize}

\subsection[Vita - Five]{\texorpdfstring{Vita -
Five\footnote{Authors Mentioned: Homer, Sophocles, Dante, Shakespeare,
  Cervantes}}{Vita - Five}}\label{vita---five5}

\begin{itemize}
\tightlist
\item
  Virgil. \emph{Virgil's Eclogues}. University of Pennsylvania Press.
  2011.\footnote{Line 6, ``Iam redit et Virgo, redeunt Saturnia regna''
    ``The virgin reappears and Saturn reigns''
    https://la.wikisource.org/wiki/Eclogae\_vel\_bucolica/Ecloga\_IV}
\item
  Dante. \emph{Divine Comedy}. Modern Library. 1950
\end{itemize}

\subsection[Vita - Six]{\texorpdfstring{Vita -
Six\footnote{Authors Mentioned: Giambattista Vico, Grimm, Frobenius,
  George Santayana, Giorgio di Santillana, Anthony Trollope}}{Vita - Six}}\label{vita---six6}

\begin{itemize}
\tightlist
\item
  Rose, Lois. \emph{Steganography}.\footnote{This is a fictional book
    created by John Crowley.}
\item
  Bates, Lizzie. \emph{Stories of the Flowers}. 1872.
\item
  Perlesvaus. \emph{The High History of the Holy Graal}. 1898.
\item
  Sprat, Thomas. \emph{The History of the Royal Society of London, for
  the Improving of Natural Knowledge}. 1667.
\item
  Frazer, James George. \emph{The Golden Bough: A Study in Magic and
  Religion}. 1932.
\item
  Longfellow, Henry Wadsworth. \emph{The Golden Legend}. 1878.
\item
  Apuleius. \emph{The Golden Ass}. Penguin UK, 1998.
\item
  Hyginus. \emph{Poeticon Astronomicon}. 1502.
\item
  Huizinga, J. \emph{The Waning of the Middle Ages: A Study of the Forms
  of Life, Thought and Art in France and The Netherlands in the XIVth
  and XVth Centuries}. Pickle Partners Publishing, 2016.
\item
  Burgess, Thornton W. \emph{Old Mother West Wind}. Courier Corporation,
  2012.\footnote{Mentioned by Crowley as ``Old Mother Westwind''}
\item
  Miss Martha's History Book\footnote{A completely unnamed history book
    belonging to Miss Martha, Pierce's uncle Sam's Houskeeper.}
\end{itemize}

\subsection[Vita - Seven]{\texorpdfstring{Vita -
Seven\footnote{Author's Mentioned: Æneas.}}{Vita - Seven}}\label{vita---seven8}

\begin{itemize}
\tightlist
\item
  Shakespeare, William. \emph{Shakespeare's King Henry VIII}. 1912.
\end{itemize}

\subsection[Lucrum - One]{\texorpdfstring{Lucrum -
One\footnote{Author's Mentioned: Aeschylus, Euripides, Leland, Polydore
  Virgil, Charles Dickens, Sir Walter Scott,}}{Lucrum - One}}\label{lucrum---one9}

\begin{itemize}
\tightlist
\item
  Sidney, Phillip, Sir. \emph{The Lady of May}. 1598.
\item
  Holinshed, Raphael. \emph{Holinshed's Chronicals: Richard II,
  1398-1400, Henry IV and Henry V.} 1923.
\item
  Monardes, Nicolás. \emph{Joyfull Newes Out of the Newe Founde Worlde}.
  Walter J. Johnson Incorporated, 1970.
\item
  Plutarch and North, Thomas, Sir. \emph{Lives. Englished}. London Dent.
  1894.
\end{itemize}

\subsection[Lucrum - Three]{\texorpdfstring{Lucrum -
Three\footnote{Author's Mentioned: Gerardus Mercator, Tycho Brahe}}{Lucrum - Three}}\label{lucrum---three10}

\begin{itemize}
\tightlist
\item
  Trithemius, Johann. \emph{Johannis Trithemii \ldots{}
  Steganographia\ldots Illustrata a Wolfgango Ernesto Heidel}. 1676.
\item
  Dee, John. \emph{Propædeumata Aphoristica Ioannis Dee, Londinensis, De
  Præstantioribus Quibusdam Naturæ Virtutibus}. 1568.
\item
  \emph{Talbot's Unnamed Enochian Manuscript}\footnote{This is a
    fictional book created by John Crowley.}
\end{itemize}

\subsection[Lucrum - Four]{\texorpdfstring{Lucrum -
Four\footnote{Authors Mentioned: Giordano Bruno, Copernicus, Galileo
  Galilei, Thomas Aquinas, Dante, Aristotle, Hermes Trismegistus,
  Cicero, Lactantius, Plato, Virgil, Pythagoras, Plotinus, Jean-François
  Champollion, E. A. Wallis Budge. Margaret Mead, Aleister Crowley,
  Martin del Rio ``Delrio,'' Johannes Kepler, René Descartes, Francis
  Bacon, Isaac Newton, Proclus, Iamblichus, Paracelsus}}{Lucrum - Four}}\label{lucrum---four12}

\begin{itemize}
\tightlist
\item
  Mead, G.R.S. \emph{Thrice-greatest Hermes}. London and Benares. 1906.
\item
  Beeson, Charles H. ``A History of Magic and Experimental Science
  During the First Thirteen Centuries of Our Era. Lynn Thorndyke.''
  \emph{Classical Philology}, vol.~20, no. 2, Apr.~1925, pp.~163--67.
  https://doi.org/10.1086/360664.
\item
  Bruno, Giordano. \emph{Opera Omnia}. 1854.
\item
  ---. \emph{Jordani Bruni Opera Latine Conscripta}. 1962.
\end{itemize}

\subsection[Lucrum - Five]{\texorpdfstring{Lucrum -
Five\footnote{Authors Mentioned: Giordano Bruno, Agrippa (unsure which
  one), Gerolamo Cardano, John Galsworthy,}}{Lucrum - Five}}\label{lucrum---five13}

\begin{itemize}
\tightlist
\item
  Kraft, Fellowes. \emph{The Court of Silk and Blood}.\footnote{This is
    a fictional book created by John Crowley.}
\item
  ---. \emph{A Passage at Arms}.\footnote{This is a fictional book
    created by John Crowley.}
\item
  Unnamed \emph{Nancy Drew Mystery Stories} Novel.
\item
  Unnamed \emph{Mr.~Moto} novel by John P. Marquand.
\end{itemize}

\subsection{Lucrum - Six}\label{lucrum---six}

\begin{itemize}
\tightlist
\item
  \emph{TV Guide}.
\item
  \emph{Faraway Crier}.\footnote{This is a fictional book created by
    John Crowley.}
\item
  Kraft, Fellowes. \emph{Sit Down, Sorrow}.\footnote{This is a fictional
    book created by John Crowley.}\footnote{A memoir of Fellowes Kraft.}
\end{itemize}

\subsection[Lucrum - Seven]{\texorpdfstring{Lucrum -
Seven\footnote{Author's Mentioned: Karl Marx}}{Lucrum - Seven}}\label{lucrum---seven15}

\begin{itemize}
\tightlist
\item
  Seneschal, Jean Le, and Gaston Raynaud. \emph{Les Cent Ballades\,:
  Poème Du XIV{[}E{]} Siècle}. 1905, p221.
\item
  \emph{The Eighteenth Brumaire of Louis Bonaparte}, Charles H. Kerr,
  Chicago, 1907.\footnote{````Hegel says somewhere that all great
    historic facts and personages recur twice. He forgot to add:''Once
    as tragedy, and again as farce.'' Caussidiere for Danton, Louis
    Blanc for Robespierre, the ``Mountain'' of 1848-51 for the.
    ``Mountain'' of 1793-95, the Nephew for the Uncle. The identical
    caricature marks also the conditions under which the second edition
    of the eighteenth Brumaire is issued.''}
\item
  \emph{Phaeton's Car}.\footnote{This is a fictional book created by
    John Crowley.}\footnote{Fictitious book apparently a play on Francis
    Bacon's \emph{Of Vicissitude of Things} ``Phaeton's car went but a
    day: and the three years drought in the time of Elias, was but
    particular, and left people alive.'' Para. 60. It is apparently a
    fictional renaming of \emph{The Chariot of the Gods} (Von Däniken,
    Erich. \emph{Chariots of the Gods?: Unsolved Mysteries of the Past}.
    1973.)}
\item
  \emph{Worlds in Division}.\footnote{This is a fictional book created
    by John Crowley.}
\item
  \emph{Dawn of the Druids}.\footnote{This is a fictional book created
    by John Crowley.}
\end{itemize}

\subsection[Lucrum - Eight]{\texorpdfstring{Lucrum -
Eight\footnote{Authors Mentioned: Rablelais}}{Lucrum - Eight}}\label{lucrum---eight19}

\begin{itemize}
\tightlist
\item
  Lanier, Sidney and Wyeth, N. C., Illustrator. \emph{The Boy's King
  Arthur. Scribner's Illustrated Classics. Sir Thomas Malory's History
  of King Arthur and His Knights of the Round Table}. Charles Scribner's
  Sons. 1937.
\item
  Briggs, Austin. \emph{Flash Gordon}. 1940-1948.
\end{itemize}

\subsection[Fratres - Two]{\texorpdfstring{Fratres -
Two\footnote{Authors Mentioned: Norton, Norris, Nofzinger, Mackenzie,
  Macauley, Macdonald, Ross Lockbridge, Joseph Lincoln, Shellabarger,
  Costain, Milton}}{Fratres - Two}}\label{fratres---two20}

\begin{itemize}
\tightlist
\item
  \emph{New York Times}.
\item
  \emph{Phaeton}'\emph{s Car}.\footnote{This is a fictional book created
    by John Crowley.}\footnote{Fictitious book apparently a play on
    Francis Bacon's \emph{Of Vicissitude of Things} ``Phaeton's car went
    but a day: and the three years drought in the time of Elias, was but
    particular, and left people alive.'' Para. 60. It is apparently a
    fictional renaming of \emph{The Chariot of the Gods} (Von Däniken,
    Erich. \emph{Chariots of the Gods?: Unsolved Mysteries of the Past}.
    1973.)}
\item
  Niblick, Helen. \emph{The Way's Far Turning}.\footnote{This is a
    fictional book created by John Crowley.}
\item
  Kraft, Fellowes. \emph{The Book of a Hundred Chapters}.\footnote{This
    is a fictional book created by John Crowley.}
\item
  \emph{Dictionary of National Biography}. 1961.
\item
  Benians, Ernest Alfred. \emph{The Cambridge Modern History}.
  Cambridge\,: UP, 1907.
\item
  Herbermann, Charles George. \emph{The Catholic Encyclopedia: An
  International Work of Reference on the Constitution, Doctrine,
  Discipline, and History of the Catholic Church}. 1908.
\item
  Barr, Franks Walker. Unnamed collected works of histories.
\item
  \emph{Unnamed Bible}.
\item
  Stedman, Joseph F. \emph{My Sunday Missal: Using New Translation From
  New Testament and a Simplified Method of Following Mass With an
  Explanation Before Each Mass of Its Theme}. Confraternity of the
  Precious Blood, 1961.\footnote{Title given in Aegypt is \emph{Our
    Sunday Missal} which could have been a misremembering of Crowley's
    part, or another slight alteration of the book's world from our own.
    Either way, this is the closest thing I can see Pierce having on his
    shelf.}
\item
  Shakespeare, William. \emph{Unknown volumes of the Yale Shakespeare}
\item
  Dee, John. \emph{Monas hieroglyphica}. Antwerp, 1564.
\item
  Thorndike, Lynn. \emph{Four Unnamed Volumes}.\footnote{While the four
    volumes are not mentioned by name, we can make some guesses as to
    which they are since they are a part of the ``secondary sources''
    mentioned for Pierce's novel. ( i.e.~\emph{The Place of Magic in the
    Intellectual History of Europe}, \emph{The History of Medieval
    Europe}, up to any four of the eight volumes of the \emph{History of
    Magic and Experimental Science}, \emph{Check-list of Rotographs in
    the History of Natural and Occult Sciences},
    \emph{\href{http://hdl.handle.net/2027/mdp.39015025028716}{\emph{The
    Sphere of Sacrobosco and its Commentators,}}}, etc. )}
\item
  Sacrobosco, Earl. \emph{Unnamed Astronomy Textbook}.\footnote{This is
    a fictional book created by John Crowley.}
\item
  Lewis, Charlton T. \emph{Lewis and Short Latin Dictionary}. 1966.
\item
  Davidson, Gustav. \emph{A DICTIONARY OF ANGELS\,: INCL. THE FALLEN
  ANGELS}. 1971.\footnote{While this is not named outright and is simply
    referred to as ``a dictionary of angels'' it is by far the most
    famous published prior to the 1980s and is likely the one Crowley
    was referring to.}
\item
  \emph{Unnamed Big Dictionary}.
\item
  \emph{Unnamed Big Picture Book of Clocks}.\footnote{There are a few
    possibilities here as to what this books could have been. I'm simply
    going to list the most likely candidates here: ``The Book of
    Clocks'' by David Thompson (1971), ``Clocks and Watches'' by Eric
    Bruton (1973), ``The Art of the Clockmaker'' by Tardy (early 1970s),
    ``The Lore of the Clock'' by G. H. Baillie (1973).}
\item
  \emph{Unknown Volumes of Encyclopedia Britannica}. 1939.\footnote{The
    Encyclopedia Britannica became the first encyclopedia to undergo
    continual revisions in 1933. Crowley mentions ``Several volumes of
    the 1939 edition'' in this chapter, but there was not a specific
    edition published that year and as such it is impossible to easily
    tell what volumes of the 1933 14th edition were reprinted during
    1939 specifically. Again this could be either intentional by Crowley
    subtly showing the differences between our world and the book's, or
    he simply could have chosen a date. Personally, I think it is likely
    the former.}
\item
  \emph{Unknown Big Shakespeare}.
\item
  \emph{Unknown Big Bible}.
\item
  Murphy, John and et. al.~\emph{The Holy Bible : translated from the
  Latin Vulgate: \ldots{}} University of Toronto, 1914.\footnote{The
    Douay-Rheims Bible, sometimes known simply as the Douai Bible was a
    translation for the Latin Vulgate completed by the English College,
    Douai in 1582 \& 1609.}
\end{itemize}

\subsection[Fratres - Three]{\texorpdfstring{Fratres -
Three\footnote{Authors Mentioned: Marsilio Ficino, Robert Fludd}}{Fratres - Three}}\label{fratres---three28}

\begin{itemize}
\tightlist
\item
  Kraft, Fellowes. \emph{Bruno's Journey}\footnote{This is a fictional
    book created by John Crowley.}\footnote{This was Kraft's first book.}
\item
  ---. \emph{Under Saturn}.
\item
  ---. \emph{Darkling Plain}.
\item
  Fludd, Robert. \emph{Unnamed book}.\footnote{This can only mean
    Fludd's \emph{Utriusque Cosmi}, however, there were no publications
    in english prior to the late 20th century. I believe it is instead
    referring to the following, but there will never be a way to tell
    for certain. (Paracelsus. \emph{The Hermetic and Alchemical Writings
    of Aureolus Philippus Theophrastus Bombast, of Hohenheim, Called
    Paracelsus the Great}. 1894.)}
\end{itemize}

\subsection{Fratres - Four}\label{fratres---four}

\begin{itemize}
\tightlist
\item
  ``\emph{Death's dream kingdom}''. Unknown Poem.\footnote{This must
    refer to T. S. Eliot's ``The Hollow Men'' from \emph{Poems,
    1909--1925} (Faber \& Faber Limited, 1925).}
\item
  \emph{Picatrix: Das Ziel des Weisen von Pseudo-Magriti}, aus dem
  Arabischen ins Deutsche übersetzt von Hellmut Ritter und Martin
  Plessner {[}\emph{Picatrix: The Goal of the Wise Man by
  Pseudo-Magriti}, translated from Arabic into German by Ritter and
  Plessner{]}. London: Warburg Institute, 1962
\item
  Kraft, Fellowes. \emph{Unfinished Final Novel}.
\end{itemize}

\subsection[Fratres - Five]{\texorpdfstring{Fratres -
Five\footnote{Authors Mentioned: Cicero, Quintilian, St.~Albertus Magus,
  St.~Thomas Aquinas, Aristotle, St.~Jerome, St.~Ambrose, Augustine of
  Hippo, Lactantias, Horapollo}}{Fratres - Five}}\label{fratres---five31}

\begin{itemize}
\tightlist
\item
  Thomas, Saint. \emph{The ``Summa Theologica'' of St.~Thomas Aquinas}.
  1917.
\item
  Magus, Saint Albertus. \emph{Commentaries on The ``Summa
  Theologica''}.
\item
  ---. \emph{On Sleeping and Being Awake''}.
\item
  ---. \emph{Book of Secrets''}.
\item
  de Sacrobosco, Johannes. \emph{Tractatus de Sphæra}.
\item
  Cecco of Ascoli. \emph{Commentary on the Sphere}.1250.
\item
  Solomon. \emph{Shadows of Ideas}.\footnote{A completely unnamed
    history book belonging to Miss Martha, Pierce's uncle Sam's
    Houskeeper.}
\item
  Ficcino, Marsilo. \emph{De vita coelitis comparanda}.
\item
  Horapollo. \emph{Hieroglyphica}. \textasciitilde1419.
\item
  Iambilichus. \emph{On the Mysteries of the Egyptians}.
\item
  Apuleius. \emph{The Golden Ass}. 158-180.
\item
  al-Ḥakīm, Ghāyat. \emph{Picatrix}.
\end{itemize}

\subsection{Fratres - Six}\label{fratres---six}

\begin{itemize}
\tightlist
\item
  \emph{Unnamed Children's Book}.\footnote{I spent way too much time
    trying to track this book down. As best I can tell, there is no
    surviving children's ``book about the family of mice who go
    traveling in a balloon'' that was published in 1978 or before.}
\end{itemize}

\subsection[Fratres - Seven]{\texorpdfstring{Fratres -
Seven\footnote{Authors Mentioned: Cecco of Ascoli, Plato, Aristotle,
  Pythagoras, Plato, Epicurus, Marsilio Ficino}}{Fratres - Seven}}\label{fratres---seven33}

\begin{itemize}
\tightlist
\item
  \emph{Psalms 87 / Hebrew - English Bible / Mechon-Mamre}.
  \url{mechon-mamre.org/p/pt/pt2687.htm}.
\item
  Ficcino, Marsilio. \emph{Pimander}. 1543.\footnote{Klutstein, Ilana.
    ``Ficino's Hermetic Translations: English Translation of His Latin
    Pimander.'' \emph{Huji}, Feb.~2017,
    www.academia.edu/31423671/Ficinos\_Hermetic\_Translations\_English\_Translation\_of\_His\_Latin\_Pimander.}
\end{itemize}

\subsection[Fratres - Eight]{\texorpdfstring{Fratres -
Eight\footnote{Authors Mentioned: Pythagoras, Zoroaster}}{Fratres - Eight}}\label{fratres---eight36}

\begin{itemize}
\tightlist
\item
  de Sacrobosco, Johannes. \emph{Tractatus de Sphæra}.
\item
  Bruno, Giordano. \emph{Ars Memoriae}. 1582.\footnote{The book itself
    is not mentioned, instead Crowley only says, ``He lectured on the
    \emph{Ars memoriae}, and let it be whispered that he had fled a
    Dominican monastery\ldots{}'' The time-line for publication is still
    some time off at this stage of the history of Bruno, but it seems
    obvious that what he would be teaching is in fact the corpus of what
    would later become \emph{Ars Memoriae}.}
\item
  Iambilichus. \emph{On the Mysteries of the Egyptians}.
\item
  Agrippa, Heinrich Cornelius. \emph{De Occulta Philosophia}. 1533
\item
  \emph{ORPHIC HYMNS 1-40 - Theoi Classical Texts Library}.
  www.theoi.com/Text/OrphicHymns1.html.
\item
  Lull, Ramón. \emph{Ars Magna}. 1501.
\item
  Copernicus, Nicolaus. De Revolutionibus Orbium Coelestium. 1543.
\item
  \emph{Das Nibelungenlied: Song of the Nibelungs}. Yale UP, 2006.
\item
  Malory, Thomas. \emph{Le Morte Darthur: The Winchester Manuscript}.
  Oxford UP, USA, 1998.
\item
  Spenser, Edmund. \emph{The Faerie Queene}. CUP Archive, 1920.
\end{itemize}

\subsection[Fratres - Nine]{\texorpdfstring{Fratres -
Nine\footnote{Authors Mentioned: Copernicus, Aristotle, Lucretius,
  Cusanus, Antoine de la Faye}}{Fratres - Nine}}\label{fratres---nine37}

\begin{itemize}
\tightlist
\item
  Bruno, Giordano. \emph{De umbris idearum}. 1582.
\end{itemize}

\subsection[Fratres -Ten]{\texorpdfstring{Fratres
-Ten\footnote{Authors Mentioned: John Dee, Giordano Bruno, Origen,
  Descartes, Peter Ramus, Francis Bacon}}{Fratres -Ten}}\label{fratres--ten38}

\begin{itemize}
\tightlist
\item
  \emph{Unknown volume of The Catholic Encyclopedia}.
\item
  De Góngora Y Argote, Luis. \emph{The Solitudes of Luis De Góngora}.
  1968.
\item
  Kraft, Fellowes. \emph{Bitten Apples}\footnote{This is a fictional
    book created by John Crowley.}
\end{itemize}

\section{Complete List of Authors Mentioned in
Ægypt}\label{complete-list-of-authors-mentioned-in-uxe6gypt}

\begin{itemize}
\tightlist
\item
  Aeschylus
\item
  Aleister Crowley
\item
  Antoine de la Faye
\item
  Aristotle
\item
  Cecco of Ascoli
\item
  Charles Dickens
\item
  Charles Godfrey Leland
\item
  Cicero
\item
  Copernicus
\item
  Thomas B. Costain
\item
  Cusanus - Nicholas of Cusa
\item
  Dante
\item
  E. A. Wallis Budge
\item
  Epicurus
\item
  Euripides
\item
  Francis Bacon
\item
  François Rabelais
\item
  Galileo Galilei
\item
  Gerardus Mercator
\item
  Giordano Bruno
\item
  Hermes Trismegistus
\item
  Homer
\item
  Iamblichus
\item
  Isaac Newton
\item
  Jean-François Champollion
\item
  Johannes Kepler
\item
  John Dee
\item
  Joseph Lincoln
\item
  Karl Marx
\item
  Kathleen Norris\footnote{I'm reasonably sure that this is Kathleen
    Norris, but there is an argument to be made that he was referring to
    Frank Norris instead.}
\item
  Lady Mary Wortley Montagu
\item
  Lactantius
\item
  Lucretius
\item
  Luis de Góngora
\item
  Mackenzie
\item
  Macdonald
\item
  Margaret Mead
\item
  Margaret Nofzinger
\item
  Marsilio Ficcino
\item
  Martin del Rio ``Delrio''
\item
  Mary Norton
\item
  Miguel de Cervantes
\item
  Origen
\item
  Paracelsus
\item
  Percy Bysshe Shelley
\item
  Peter Ramus
\item
  Plato
\item
  Plotinus
\item
  Polydore Virgil
\item
  Proclus
\item
  Pythagoras
\item
  René Descartes
\item
  Robert Fludd
\item
  Ross Lockbridge
\item
  Samuel Shellabarger
\item
  Sir Walter Scott
\item
  Sophocles
\item
  Thomas Aquinas
\item
  Thomas Babington Macaulay
\item
  Thomas Milton
\item
  Tycho Brahe
\item
  Virgil
\item
  William Shakespeare
\item
  Zoroaster
\end{itemize}
